% GENERATED FILE. Edit canonical lessons, then run npm run generate:books.
% canonical-source-hash: fnv1a64:1316d5b7b64ece46
% canonical-lessons: MR-C377-prabhav, MR-C377-dabav, MR-C377-protsahan-dene, MR-C377-prerna, MR-C377-khatakne

\chapter{Influence and Pressure}
\label{ch:mr-influence-and-pressure}
\hlchaptermodality{377}

\begin{chapteropening}
\textbf{By the end of this chapter:} \emph{I can talk about influence, pressure and encouragement.}

Complete the last lesson of chapter 377: I can talk about influence, pressure and encouragement.
\end{chapteropening}

\section[\texorpdfstring{\mr{प्रभाव} (prabhāv)}{प्रभाव (prabhāv)}]{\mr{प्रभाव} (prabhāv) — an influence}

\label{lesson:MR-C377-prabhav}



\begin{quote}
Before the new one: say the Marathi for quality, then the Marathi for a standard.
\end{quote}

\subsection*{You'll want to know: \mr{प्रभाव}}
\textbf{\mr{प्रभाव}} — \emph{prabhāv} — \textquotedblleft{}an influence\textquotedblright{}.

\begin{quote}
\mr{त्याच्यावर} \mr{वडिलांचा} \mr{मोठा} \mr{प्रभाव} \mr{आहे}.
\end{quote}

\begin{quote}
\emph{tyācyāvar vaḍilāñcā moṭhā prabhāv āhe.}
\end{quote}

\begin{quote}
His father has a great influence on him.
\end{quote}

Masculine.

X \mr{वर} Y \mr{चा} \mr{प्रभाव}, Y's influence on X.

\begin{cousinweb}
\textbf{\mr{प्रभाव}} (\emph{prabhāv}) is a tatsama, taken from Sanskrit \emph{prabhāva}, \textquotedblleft{}power, might\textquotedblright{} (Monier-Williams).
\end{cousinweb}

\begin{grammarlens}[title={What you've built}]
One more word for giving reasons and opinions.
\end{grammarlens}

\subsection*{Your turn}
\begin{itemize}
\raggedright
  \item \emph{Say it:} \emph{prabhāv}
  \item \emph{Say it:} \emph{prabhāv}, once more
  \item \emph{Say it:} \emph{tyācyāvar vaḍilāñcā moṭhā prabhāv āhe.} — His father has a great influence on him.
  \item \emph{Recall:} say the Marathi for quality, then the Marathi for a standard, then say \emph{prabhāv} again
\end{itemize}

\begin{tcolorbox}[breakable,colback=teal!4,colframe=teal!35!black,title={Before you move on}]
What does \mr{प्रभाव} mean? (\textquotedblleft{}An influence\textquotedblright{}.) Say it once more.
\end{tcolorbox}
\section[\texorpdfstring{\mr{दबाव} (dabāv)}{दबाव (dabāv)}]{\mr{दबाव} (dabāv) — outside pressure; pressure}

\label{lesson:MR-C377-dabav}



\begin{quote}
Before the new one: say the Marathi for a standard, then the Marathi for an influence.
\end{quote}

\subsection*{You'll want to know: \mr{दबाव}}
\textbf{\mr{दबाव}} — \emph{dabāv} — \textquotedblleft{}outside pressure; pressure\textquotedblright{}.

\begin{quote}
\mr{घरच्यांच्या} \mr{दबावामुळे} \mr{तिने} \mr{नोकरी} \mr{सोडली}.
\end{quote}

\begin{quote}
\emph{gharacyāñcyā dabāvāmuḷe tine nokrī soḍlī.}
\end{quote}

\begin{quote}
Under pressure from her family, she left her job.
\end{quote}

Masculine.

\mr{दबाव} \mr{आणणे}, to put pressure on.

\begin{cousinweb}
\textbf{\mr{दबाव}} (\emph{dabāv}) is built on \mr{दाबणे}, \textquotedblleft{}to press\textquotedblright{}, with the noun ending -\mr{आव}; CDIAL lists Marathi \emph{dābṇẽ}, \textquotedblleft{}to press\textquotedblright{}, under \emph{*dabb-}, \textquotedblleft{}press\textquotedblright{} (CDIAL 6173).
\end{cousinweb}

\begin{grammarlens}[title={What you've built}]
One more word for giving reasons and opinions.
\end{grammarlens}

\subsection*{Your turn}
\begin{itemize}
\raggedright
  \item \emph{Say it:} \emph{dabāv}
  \item \emph{Say it:} \emph{dabāv}, once more
  \item \emph{Say it:} \emph{gharacyāñcyā dabāvāmuḷe tine nokrī soḍlī.} — Under pressure from her family, she left her job.
  \item \emph{Recall:} say the Marathi for a standard, then the Marathi for an influence, then say \emph{dabāv} again
\end{itemize}

\begin{tcolorbox}[breakable,colback=teal!4,colframe=teal!35!black,title={Before you move on}]
What does \mr{दबाव} mean? (\textquotedblleft{}Outside pressure; pressure\textquotedblright{}.) Say it once more.
\end{tcolorbox}
\section[\texorpdfstring{\mr{प्रोत्साहन} \mr{देणे} (protsāhan deṇe)}{प्रोत्साहन देणे (protsāhan deṇe)}]{\mr{प्रोत्साहन} \mr{देणे} (protsāhan deṇe) — to encourage}

\label{lesson:MR-C377-protsahan-dene}



\begin{quote}
Before the new one: say the Marathi for an influence, then the Marathi for outside pressure.
\end{quote}

\subsection*{You'll want to know: \mr{प्रोत्साहन} \mr{देणे}}
\textbf{\mr{प्रोत्साहन} \mr{देणे}} — \emph{protsāhan deṇe} — \textquotedblleft{}to encourage\textquotedblright{}.

\begin{quote}
\mr{शिक्षकांनी} \mr{मला} \mr{गाणं} \mr{शिकायला} \mr{प्रोत्साहन} \mr{दिलं}.
\end{quote}

\begin{quote}
\emph{śikṣakānnī malā gāṇaṁ śikāylā protsāhan dilaṁ.}
\end{quote}

\begin{quote}
The teacher encouraged me to learn singing.
\end{quote}

Past, with -\mr{ने} (he …): \mr{त्याने} \mr{प्रोत्साहन} \mr{दिलं}, \emph{tyāne protsāhan dilaṁ}.

The one encouraged takes -\mr{ला}, the aim -\mr{ला} or -\mr{साठी}.

\begin{cousinweb}
\textbf{\mr{प्रोत्साहन} \mr{देणे}} (\emph{protsāhan deṇe}) joins the tatsama \mr{प्रोत्साहन}, Sanskrit \emph{protsāhana}, \textquotedblleft{}encouraging, inciting\textquotedblright{} (Monier-Williams), to \mr{देणे}, \textquotedblleft{}to give\textquotedblright{}.
\end{cousinweb}

\begin{grammarlens}[title={What you've built}]
One more word for giving reasons and opinions.
\end{grammarlens}

\subsection*{Your turn}
\begin{itemize}
\raggedright
  \item \emph{Say it:} \emph{protsāhan deṇe}
  \item \emph{Say it:} \emph{protsāhan deṇe}, once more
  \item \emph{Say it:} \emph{śikṣakānnī malā gāṇaṁ śikāylā protsāhan dilaṁ.} — The teacher encouraged me to learn singing.
  \item \emph{Recall:} say the Marathi for an influence, then the Marathi for outside pressure, then say \emph{protsāhan deṇe} again
\end{itemize}

\begin{tcolorbox}[breakable,colback=teal!4,colframe=teal!35!black,title={Before you move on}]
What does \mr{प्रोत्साहन} \mr{देणे} mean? (\textquotedblleft{}To encourage\textquotedblright{}.) Say it once more.
\end{tcolorbox}
\section[\texorpdfstring{\mr{प्रेरणा} (prerṇā)}{प्रेरणा (prerṇā)}]{\mr{प्रेरणा} (prerṇā) — inspiration; motivation}

\label{lesson:MR-C377-prerna}



\begin{quote}
Before the new one: say the Marathi for outside pressure, then the Marathi for to encourage.
\end{quote}

\subsection*{You'll want to know: \mr{प्रेरणा}}
\textbf{\mr{प्रेरणा}} — \emph{prerṇā} — \textquotedblleft{}inspiration; motivation\textquotedblright{}.

\begin{quote}
\mr{आईच्या} \mr{कष्टांतून} \mr{मला} \mr{प्रेरणा} \mr{मिळाली}.
\end{quote}

\begin{quote}
\emph{āīcyā kaṣṭāntūn malā prerṇā miḷālī.}
\end{quote}

\begin{quote}
I drew inspiration from my mother's hard work.
\end{quote}

Feminine.

X \mr{ची} \mr{प्रेरणा} \mr{घेणे}, to take inspiration from X.

\begin{cousinweb}
\textbf{\mr{प्रेरणा}} (\emph{prerṇā}) is a tatsama, taken from Sanskrit \emph{preraṇa}, \textquotedblleft{}driving on, urging\textquotedblright{} (Monier-Williams).
\end{cousinweb}

\begin{grammarlens}[title={What you've built}]
One more word for giving reasons and opinions.
\end{grammarlens}

\subsection*{Your turn}
\begin{itemize}
\raggedright
  \item \emph{Say it:} \emph{prerṇā}
  \item \emph{Say it:} \emph{prerṇā}, once more
  \item \emph{Say it:} \emph{āīcyā kaṣṭāntūn malā prerṇā miḷālī.} — I drew inspiration from my mother's hard work.
  \item \emph{Recall:} say the Marathi for outside pressure, then the Marathi for to encourage, then say \emph{prerṇā} again
\end{itemize}

\begin{tcolorbox}[breakable,colback=teal!4,colframe=teal!35!black,title={Before you move on}]
What does \mr{प्रेरणा} mean? (\textquotedblleft{}Inspiration; motivation\textquotedblright{}.) Say it once more.
\end{tcolorbox}
\section[\texorpdfstring{\mr{खटकणे} (khaṭakṇe)}{खटकणे (khaṭakṇe)}]{\mr{खटकणे} (khaṭakṇe) — to bother; to jar on}

\label{lesson:MR-C377-khatakne}



\begin{quote}
Before the new one: say the Marathi for to encourage, then the Marathi for inspiration.
\end{quote}

\subsection*{You'll want to know: \mr{खटकणे}}
\textbf{\mr{खटकणे}} — \emph{khaṭakṇe} — \textquotedblleft{}to bother; to jar on\textquotedblright{}.

\begin{quote}
\mr{त्याचं} \mr{उद्धट} \mr{बोलणं} \mr{मला} \mr{खटकलं}.
\end{quote}

\begin{quote}
\emph{tyācaṁ uddhaṭ bolṇaṁ malā khaṭaklaṁ.}
\end{quote}

\begin{quote}
His rude way of talking bothered me.
\end{quote}

Past (to him …): \mr{त्याला} \mr{खटकलं}, \emph{tyālā khaṭaklaṁ}.

Dative: \mr{मला} \mr{खटकलं}, it bothered me; what jars is the subject.

\begin{cousinweb}
\textbf{\mr{खटकणे}} (\emph{khaṭakṇe}) is a Marathi verb whose origin is not settled.
\end{cousinweb}

\begin{grammarlens}[title={What you've built}]
One more word for giving reasons and opinions.
\end{grammarlens}

\subsection*{Your turn}
\begin{itemize}
\raggedright
  \item \emph{Say it:} \emph{khaṭakṇe}
  \item \emph{Say it:} \emph{khaṭakṇe}, once more
  \item \emph{Say it:} \emph{tyācaṁ uddhaṭ bolṇaṁ malā khaṭaklaṁ.} — His rude way of talking bothered me.
  \item \emph{Recall:} say the Marathi for to encourage, then the Marathi for inspiration, then say \emph{khaṭakṇe} again
\end{itemize}

\begin{tcolorbox}[breakable,colback=teal!4,colframe=teal!35!black,title={Before you move on}]
What does \mr{खटकणे} mean? (\textquotedblleft{}To bother; to jar on\textquotedblright{}.) Say it once more.
\end{tcolorbox}
